\subsection{Discussion of Metrics Comparison}
\label{sec:metrics_comparison}

In addition to runtime-based metrics, we incorporate text-based metrics to assess LLMs' performance in the test generation task on the $Methods2Test_{filter}$ dataset. Notably, we include a syntax error rate, a measure that provides valuable insights without requiring code execution.

\begin{table}[htbp]
    \centering
    \small
    \caption{Comparison of LLMs across different metrics on Defects4J and $Methods2Test_{filter}$ benchmarks.}
    \label{tab:metrics_compirason}
    \resizebox{\textwidth}{!}{
    \begin{tabular}{lcccccccccc}
    \toprule
    \multirow{2}{*}{\textbf{LLMs}}
    & \multicolumn{5}{c}{\textbf{Defects4J}} & \multicolumn{4}{c}{\textbf{Methods2Test$_{\textbf{filter}}$}} \\
    \cmidrule(lr){2-6} \cmidrule(lr){7-10}
     & \textbf{Correct} & \textbf{Passing} & \textbf{Failing} & \textbf{Build Error} & \textbf{Syntax Error} & \textbf{EM} & \textbf{BLEU} & \textbf{CodeBLEU} & \textbf{Syntax Error} \\

    \midrule
    CodeBERT & 5.56\% & 8.41\% & 5.99\% & 78.90\% & 6.70\% & 0.21\% & 11.96\% & 18.23\% & 9.57\% \\
    GraphCodeBERT & 11.77\% & 14.02\% & 11.19\% & 65.89\% & 8.89\% & 0.55\% & 12.65\% & 18.51\% & 10.99\% \\
    UniXcoder & 6.42\% & 8.12\% & 10.48\% & 69.02\% & 12.37\% & 0.46\% & 12.41\% & 18.82\% & 10.23\% \\
    CodeT5-base & 15.91\% & 18.76\% & 14.90\% & 61.36\% & 4.98\% & 0.29\% & 6.11\% & 23.89\% & 10.09\% \\
    CodeT5p-220m & 17.75\% & 19.79\% & 27.07\% & 50.26\% & 2.88\% & 0.30\% & 12.11\% & 23.00\% & 3.73\% \\
    PLBART-large & 17.50\% & 20.25\% & 20.63\% & 52.96\% & 6.16\% & 0.05\% & 9.43\% & 22.16\% & 6.47\% \\
    CodeGPT & 7.66\% & 9.64\% & 9.81\% & 39.67\% & 40.88\% & 0.09\% & 8.28\% & 21.50\% & 15.47\% \\
    StarCoder2-7b & 20.78\% & 23.10\% & 23.16\% & 41.45\% & 12.30\% & 0.27\% & 5.56\% & 22.34\% & 5.92\% \\
    CodeLlama-7b & 29.50\% & 32.14\% & 7.86\% & 30.87\% & 29.13\% & 0.37\% & 9.26\% & 25.11\% & 4.96\% \\
    DeepSeek-Coder-6b & 33.68\% & 36.02\% & 9.00\% & 21.36\% & 33.62\% & 0.43\% & 8.61\% & 24.96\% & 5.18\% \\
    \bottomrule
    \end{tabular}
    }
\end{table}


As presented in Table~\ref{tab:metrics_compirason}, all LLMs perform significantly worse on text-based metrics, with EM scores ranging from only 0.04\% to 0.55\%, a stark contrast to their performance on runtime-based metrics. 
In terms of BLEU and CodeBLEU scores, LLMs still underperform, with a BLEU score range from 5.56\% to 12.65\% and a CodeBLEU score range from 18.23\% to 25.11\%, remaining unexpectedly low compared to their performance on the other two tasks.
Despite the low EM scores, the generated code maintains a relatively low syntax error rate, ranging from 4.96\% to 15.47\%, similar to results observed on Defects4J. Interestingly, decoder-only models also perform well in syntax error rate, likely due to the Methods2Test dataset’s filtering process, which removes overly long inputs.

We attribute this lack of correlation to the diversity of the test cases. As mentioned in Section~\ref{sec:dataset}, half of the focal methods in the Methods2Test dataset include at least two related test cases, reflecting real-world development practices. For example, the focal method \texttt{`public Collection<Attribute> encode();'} in the Methods2Test dataset has eight associated test cases, each testing the method across different specific encoding types and verifying that the correct collection is returned. 
Although we remove duplicate test cases, the overall distribution remains unchanged. This may explain why LLMs perform poorly on text-based metrics yet relatively well on runtime-based metrics: the generated test cases may test the focal method from perspectives different from those in the ground truth.
These findings indicate that, given the diversity of test cases and the real-world development context, runtime-based metrics are more suitable than text-based metrics for evaluating the effectiveness of automated test generation approaches.


\summary{
LLMs may perform poorly on text-based metrics due to the diversity of test cases. Considering realistic development scenarios, it is recommended to use runtime-based metrics to evaluate the effectiveness of the automated test generation approaches.
}
